\documentclass[11pt]{article}

\usepackage[margin=1.45in]{geometry}
\usepackage[T1]{fontenc}
\usepackage[utf8]{inputenc}
\usepackage{amsmath,amssymb}
\usepackage{enumitem}
\usepackage[dvipsnames]{xcolor}
\usepackage[colorlinks=true,linkcolor=red,urlcolor=red,citecolor=red]{hyperref}

\setlist[itemize]{itemsep=2pt,topsep=3pt}

\title{\textbf{Deep Learning in Finance}}
\author{Oualid Missaoui}
\date{}

\begin{document}

\maketitle

\section*{Course Description}

This course consists of seven sessions, each lasting 2.5 hours, and follows the
format of an advanced topics seminar. It provides an in-depth exploration of
advanced deep learning techniques with applications to financial time series
data.

Prior familiarity with deep neural networks and backpropagation derivation is
required. The course begins with an overview of the theoretical foundations
that justify deep learning models, including the virtue of complexity in return
prediction as the finance-facing instance of the overparameterization story. It
then turns to Transformers, foundation models for tabular data, and deep factor
models, where the cross-section of returns is treated as a latent factor problem
and autoencoders, instrumented PCA, and variational methods are brought to bear
on it. The remainder of the course covers GANs, Neural ODEs and neural
point processes, and Diffusion Models. Throughout, each architecture is paired
with the asset pricing literature that uses it.

\section*{Evaluation}

\begin{itemize}
  \item \textbf{Research Paper Presentation}: 40\%
  \item \textbf{Final Project}: 60\%
  \item \textbf{Home works}: optional, not graded
\end{itemize}

Weekly homeworks are provided as practice and are strongly recommended as
preparation for the final project, but they are not collected or graded.

\section*{Materials}

\begin{itemize}
  \item \href{https://github.com/omroot/DLinFinance-BaruchMFE2026/tree/main/lecture_slides}{\textcolor{red}{\textbf{Lecture slides}}}
  \item \href{https://github.com/omroot/DLinFinance-BaruchMFE2026/tree/main}{\textcolor{red}{\textbf{Companion GitHub Repository}}}: Python/PyTorch
\end{itemize}

\section*{Textbooks}

\subsection*{Primary Text}

Chris Bishop, Hugh Bishop (2023) \href{https://www.bishopbook.com/}{\textcolor{red}{\textit{Deep Learning}}}, The MIT Press.

\subsection*{Optional Texts}

\begin{itemize}
  \item Aston Zhang, Zachary C. Lipton, Mu Li, Alexander J. Smola (2023), \href{https://d2l.ai/}{\textcolor{red}{\textit{Dive into Deep Learning}}}, Cambridge University Press.
  \item Simon Prince, (2023) \href{https://udlbook.github.io/udlbook/}{\textcolor{red}{\textit{Understanding Deep Learning}}}, The MIT Press.
  \item Jakub M. Tomczak (2022) \href{https://link.springer.com/book/10.1007/978-3-030-93158-2}{\textcolor{red}{\textit{Deep Generative Modeling}}}, Springer.
  \item Ovidiu Calin, (2020) \href{https://link.springer.com/book/10.1007/978-3-030-36721-3}{\textcolor{red}{\textit{Deep Learning Architectures}}}, Springer.
  \item Daniel A. Roberts, Sho Yaida, Boris Hanin (2022) \href{https://deeplearningtheory.com/}{\textcolor{red}{\textit{The Principles of Deep Learning Theory: An Effective Theory Approach to Understanding Neural Networks}}}, Cambridge University Press.
  \item Jong Chul Ye (2022) \href{https://link.springer.com/book/10.1007/978-981-16-6046-7}{\textcolor{red}{\textit{Geometry of Deep Learning: A Signal Processing Perspective}}}, Springer.
  \item Agostino Capponi, Charles-Albert Lehalle (eds.) (2023) \href{https://www.cambridge.org/core/books/machine-learning-and-data-sciences-for-financial-markets/8BB31611662A96D0AB93A8A26E2D0D0A}{\textcolor{red}{\textit{Machine Learning and Data Sciences for Financial Markets: A Guide to Contemporary Practices}}}, Cambridge University Press.
\end{itemize}

\section*{Lectures}

\begin{itemize}

%% ---------------------------------------------------------------- Lecture 1
\item \textbf{Lecture 1: Deep Learning - Theoretical Justifications}
  \begin{itemize}
    \item \textbf{Scope:} This lecture provides a high-level theoretical perspective on
      why deep learning works, focusing on approximation, optimization, and
      generalization. Topics include the universal approximation theorem, double
      descent regimes, and overparameterization. The lecture emphasizes the current
      limitations in theoretical understanding that explain the empirical success of
      deep learning models. We also explore semi-empirical and practical methodology
      to check the generalizability of deep learning models. The lecture closes with
      the \textit{virtue of complexity} literature, which carries the benign
      overfitting and double descent arguments into return prediction and shows that
      the ridgeless limit of a heavily overparameterized forecaster can dominate its
      parsimonious counterpart out of sample.
    \item \textbf{Readings:}
      \begin{itemize}
        \item Belkin, M., Hsu, D., Ma, S., and Mandal, S. (2019) \href{https://arxiv.org/pdf/1812.11118}{\textcolor{red}{Reconciling Modern Machine Learning and the Bias-Variance Tradeoff}}. Proceedings of the National Academy of Sciences (PNAS), 116(32), 15849--15854.
        \item Tomaso Poggio, Andrzej Banburski, Qianli Liao (2019) \href{https://arxiv.org/pdf/1908.09375}{\textcolor{red}{Theoretical Issues in Deep Networks: Approximation, Optimization and Generalization}} Proceedings of the National Academy of Sciences (PNAS), 117 (48) 30039-30045.
        \item Yasaman Bahria, Ethan Dyera, Jared Kaplanb, Jaehoon , (2024) \href{https://www.pnas.org/doi/epub/10.1073/pnas.2311878121}{\textcolor{red}{Explaining neural scaling laws}} Proceedings of the National Academy of Sciences (PNAS), 117 (48) 30039-30045.
        \item Charles Martin and Michael Mahoney (2021), \href{https://arxiv.org/pdf/1810.01075}{\textcolor{red}{Implicit self-regularization in deep neural networks}}
        \item Oualid Missaoui (2025), \href{https://papers.ssrn.com/sol3/papers.cfm?abstract_id=5529459}{\textcolor{red}{Inductive Triplet Fine Tuning for Small Language Models}}
        \item Bryan Kelly et al. (2023) \href{https://economics.yale.edu/sites/default/files/2024-01/The\%20Journal\%20of\%20Finance\%20-\%202023\%20-\%20KELLY\%20-\%20The\%20Virtue\%20of\%20Complexity\%20in\%20Return\%20Prediction\%20\%281\%29.pdf}{\textcolor{red}{The Virtue of Complexity in Return Prediction}}
        \item Bryan Kelly et al. (2025) \href{https://papers.ssrn.com/sol3/papers.cfm?abstract_id=5346842}{\textcolor{red}{Understanding The Virtue of Complexity}}
        \item Bryan Kelly et al. (2024) \href{https://arxiv.org/pdf/2402.06635}{\textcolor{red}{Large (and Deep) Factor Models}}
      \end{itemize}
  \end{itemize}

%% ---------------------------------------------------------------- Lecture 2
\item \textbf{Lecture 2: Transformers for Financial Time Series Forecasting}
  \begin{itemize}
    \item \textbf{Scope:} This lecture introduces the Transformers architecture and
      discusses its advantages over recurrent based sequence to sequence
      architectures. It also illustrates its application to financial time series
      forecasting. The lecture then turns to deep learning statistical arbitrage,
      where attention and convolutional sequence models are used to map the residual
      time series of an arbitrage portfolio into a trading policy, as a full-scale
      application of sequence modeling to a cross-sectional trading problem.
    \item \textbf{Readings:}
      \begin{itemize}
        \item Chapter 12, Transformers, from the primary textbook.
        \item Vaswani et al, (2017) \href{https://arxiv.org/pdf/1706.03762}{\textcolor{red}{Attention is All you Need}} NIPS.
        \item Edmond Lezmi, Jiali Xu, (2023) \href{https://papers.ssrn.com/sol3/papers.cfm?abstract_id=4375798}{\textcolor{red}{Time Series Forecasting with Transformer Models and Application to Asset Management}}, Amundi Institute.
        \item Jorge Guijarro-Ordonez et al. (2022), \href{https://arxiv.org/pdf/2106.04028}{\textcolor{red}{Deep Learning Statistical Arbitrage}}
      \end{itemize}
  \end{itemize}

%% ---------------------------------------------------------------- Lecture 3
\item \textbf{Lecture 3: Foundation Models for Tabular Data}
  \begin{itemize}
    \item \textbf{Scope:} This lecture asks what a foundation model means for tabular
      data, the one regime where gradient boosted trees have long resisted deep
      learning. It develops the Bayesian view of supervised learning that underpins
      \textit{prior-data fitted networks} (PFNs): a transformer is pre-trained on
      synthetic datasets drawn from a prior over data generating mechanisms, and at
      inference time it approximates the posterior predictive distribution in a
      single forward pass, with no gradient updates and no per-dataset training. We
      cover the PFN training principle, the TabPFN architecture and its attention
      over both rows and features, and the connection to in-context learning in
      large language models. We then follow the scaling line from TabPFN v2 through
      TabICL and TabFM, and close on the asset pricing application: FinPFN and
      meta-learning for return prediction under shifting market regimes.
    \item \textbf{Readings:}
      \begin{itemize}
        \item Samuel M\"uller, Noah Hollmann, Sebastian Pineda, Josif Grabocka, Frank Hutter (2022) \href{https://arxiv.org/abs/2112.10510}{\textcolor{red}{Transformers Can Do Bayesian Inference}}, ICLR
        \item Noah Hollmann, Samuel M\"uller, Katharina Eggensperger, Frank Hutter (2023) \href{https://arxiv.org/abs/2207.01848}{\textcolor{red}{TabPFN: A Transformer That Solves Small Tabular Classification Problems in a Second}}, ICLR
        \item Noah Hollmann et al. (2025) \href{https://www.nature.com/articles/s41586-024-08328-6}{\textcolor{red}{Accurate predictions on small data with a tabular foundation model}}, Nature
        \item Jingang Qu, David Holzm\"uller, Ga\"el Varoquaux, Marine Le Morvan (2025) \href{https://arxiv.org/abs/2502.05564}{\textcolor{red}{TabICL: A Tabular Foundation Model for In-Context Learning on Large Data}}, ICML
        \item Prior Labs (2025) \href{https://arxiv.org/abs/2511.08667}{\textcolor{red}{TabPFN-2.5: Advancing the State of the Art in Tabular Foundation Models}}
        \item Shi Bin Hoo, Samuel M\"uller, David Salinas, Frank Hutter (2025) \href{https://arxiv.org/abs/2501.02945}{\textcolor{red}{From Tables to Time: Extending TabPFN-v2 to Time Series Forecasting}}
        \item Sang Michael Xie, Aditi Raghunathan, Percy Liang, Tengyu Ma (2022) \href{https://arxiv.org/abs/2111.02080}{\textcolor{red}{An Explanation of In-context Learning as Implicit Bayesian Inference}}, ICLR
        \item L\'eo Grinsztajn, Edouard Oyallon, Ga\"el Varoquaux (2022) \href{https://arxiv.org/abs/2207.08815}{\textcolor{red}{Why do tree-based models still outperform deep learning on tabular data?}}, NeurIPS
        \item Yicheng Wang, Sandro Lera (2025) \href{https://doi.org/10.1016/j.finmar.2025.101042}{\textcolor{red}{Meta-learning for return prediction in shifting market regimes}}, Journal of Financial Markets, article 101042
        \item Bryan Kelly, Boris Kuznetsov, Semyon Malamud, Teng Andrea Xu (2025) \href{https://www.nber.org/papers/w33351}{\textcolor{red}{Artificial Intelligence Asset Pricing Models}}, NBER Working Paper 33351
        \item Shihao Gu, Bryan Kelly, Dacheng Xiu (2020) \href{https://doi.org/10.1093/rfs/hhaa009}{\textcolor{red}{Empirical Asset Pricing via Machine Learning}}, Review of Financial Studies 33(5)
      \end{itemize}
  \end{itemize}

%% ---------------------------------------------------------------- Lecture 4
\item \textbf{Lecture 4: Deep Factor Models}
  \begin{itemize}
    \item \textbf{Scope:} This lecture treats the cross-section of returns as a latent
      factor problem and asks what deep learning adds to it. We start from the classical linear factor model and static PCA, together
      with the identification questions that come with it --- how many factors, and
      the Bai--Ng and eigenvalue-ratio tests that answer it --- and the
      Eckart--Young result linking PCA to low-rank approximation. We then show that
      a linear autoencoder recovers the principal subspace (Baldi--Hornik), which
      makes the autoencoder the natural nonlinear generalization of PCA.
      Instrumented PCA lets the loadings depend on observable characteristics; the
      conditional autoencoder of Gu, Kelly and Xiu makes both the loadings and the
      factors nonlinear functions of those characteristics. We review the empirical
      and simulation evidence, and close on probabilistic factors: the variational
      autoencoder and FactorVAE, which return a predictive distribution over the
      cross-section rather than a point estimate.
    \item \textbf{Readings:}
      \begin{itemize}
        \item Chapter 19, Autoencoders, from the primary textbook.
        \item Shihao Gu, Bryan Kelly, Dacheng Xiu (2021) \href{https://papers.ssrn.com/sol3/papers.cfm?abstract_id=3335536}{\textcolor{red}{Autoencoder Asset Pricing Models}}, Journal of Econometrics 222(1), 429--450.
        \item Bryan Kelly, Seth Pruitt, Yinan Su (2019) \href{https://www.nber.org/papers/w24540}{\textcolor{red}{Characteristics Are Covariances: A Unified Model of Risk and Return}}, Journal of Financial Economics 134(3).
        \item Bryan Kelly, Seth Pruitt, Yinan Su (2020) \href{https://doi.org/10.2139/ssrn.2983919}{\textcolor{red}{Instrumented Principal Component Analysis}}, SSRN Working Paper 2983919.
        \item Shihao Gu, Bryan Kelly, Dacheng Xiu (2020) \href{https://doi.org/10.1093/rfs/hhaa009}{\textcolor{red}{Empirical Asset Pricing via Machine Learning}}, Review of Financial Studies 33(5).
        \item Jushan Bai, Serena Ng (2002) \href{https://doi.org/10.1111/1468-0262.00273}{\textcolor{red}{Determining the Number of Factors in Approximate Factor Models}}, Econometrica 70(1).
        \item Pierre Baldi, Kurt Hornik (1989) \href{https://doi.org/10.1016/0893-6080\%2889\%2990014-2}{\textcolor{red}{Neural Networks and Principal Component Analysis: Learning from Examples Without Local Minima}}, Neural Networks 2(1).
        \item Diederik P. Kingma, Max Welling (2014) \href{https://arxiv.org/abs/1312.6114}{\textcolor{red}{Auto-Encoding Variational Bayes}}, ICLR.
        \item Yitong Duan, Lei Wang, Qizhong Zhang, Jian Li (2022) \href{https://ojs.aaai.org/index.php/AAAI/article/view/20369}{\textcolor{red}{FactorVAE: A Probabilistic Dynamic Factor Model Based on Variational Autoencoder for Predicting Cross-Sectional Stock Returns}}, AAAI.
        \item Guanhao Feng, Nicholas Polson, Jianeng Xu (2023) \href{https://arxiv.org/pdf/1805.01104}{\textcolor{red}{Deep Learning in Characteristics-Sorted Factor Models}}
        \item Guanhao Feng et al. (2018) \href{https://acfr.aut.ac.nz/__data/assets/pdf_file/0003/190641/FengPolsonXu_2018.pdf}{\textcolor{red}{Deep Learning Factor Alpha}}
        \item Matthew Dixon, Nicholas Polson (2023) \href{https://www.cambridge.org/core/books/abs/machine-learning-and-data-sciences-for-financial-markets/bayesian-deep-fundamental-factor-models/DEE56793B55116E1D68BB72848C7D7B4}{\textcolor{red}{Bayesian Deep Fundamental Factor Models}}, Chapter 32 in A. Capponi and C.-A. Lehalle (eds.), \textit{Machine Learning and Data Sciences for Financial Markets}, Cambridge University Press.
        \item Matthew F. Dixon, Nicholas G. Polson (2020) \href{https://arxiv.org/pdf/1903.07677}{\textcolor{red}{Deep Fundamental Factor Models}}, SIAM Journal on Financial Mathematics.
      \end{itemize}
  \end{itemize}

%% ---------------------------------------------------------------- Lecture 5
\item \textbf{Lecture 5: Generative Adversarial Models (GANs) and Adversarial Estimation in Finance}
  \begin{itemize}
    \item \textbf{Scope:} This lecture introduces GANs and presents an overview of its
      components and architecture. It also derives the underlying updating equations
      for its training. It then illustrates the applications of GANs in financial
      synthetic data generation. Finally, the lecture shows that the adversarial
      minimax structure is not confined to sample generation: in deep learning for
      asset pricing, one network parameterizes the stochastic discount factor while
      an adversarial network searches over the moment conditions that the SDF fails
      to price, which is exactly the GAN objective applied to a no-arbitrage
      estimation problem.
    \item \textbf{Readings:}
      \begin{itemize}
        \item Chapter 17, Generative Adversarial Models, from the primary textbook.
        \item Ian J. Goodfellow et al. (2014) \href{https://arxiv.org/pdf/1406.2661}{\textcolor{red}{Generative Adversarial Nets}}
        \item Adil Rengin Cetingoz, and Charles Albert Lehalle, (2025) \href{https://arxiv.org/pdf/2501.03993}{\textcolor{red}{Synthetic Data for Portfolios: A Throw of the Dice Will Never Abolish Chance}}arXiv, arXiv:2501.03993
        \item Luyang Chen et al. (2021) \href{https://arxiv.org/pdf/1904.00745}{\textcolor{red}{Deep Learning in Asset Pricing}}
      \end{itemize}
  \end{itemize}

%% ---------------------------------------------------------------- Lecture 6
\item \textbf{Lecture 6: Neural Ordinary Differential Equations (NODE) and Neural Point Processes for Financial Time Series}
  \begin{itemize}
    \item \textbf{Scope:} This lecture introduces Neural ODEs and its formulation and
      how they extend the concept of ResNets to continuous layers. It also provides
      and explanation of the architecture and implementation of NODEs in practice,
      and illustrates its application to financial time series forecasting. The
      second half extends the continuous-time latent state to event data: the
      \textit{neural Hawkes process} replaces the fixed exponential kernel of a
      classical Hawkes process with a continuous-time LSTM, and neural jump SDEs
      place the same idea directly on top of the Neural ODE machinery. We apply the
      resulting conditional intensity to a hazard problem in quantitative finance,
      modeling the arrival intensity of credit events and of self-exciting order flow.
    \item \textbf{Readings:}
      \begin{itemize}
        \item Chapter 1 and 2 from the PhD dissertatin of Patrick Kidger, \href{https://arxiv.org/pdf/2202.02435}{\textcolor{red}{On Neural Differential Equations}}, 2022.
        \item Ricky T.Q. Chen, Yulia Rubanova, Jesse Bettencourt, David Duvenaud, (2018) \href{https://proceedings.neurips.cc/paper_files/paper/2018/file/69386f6bb1dfed68692a24c8686939b9-Paper.pdf}{\textcolor{red}{Neural Ordinary Differential Equations}}, NIPS
        \item Jingsui Li; Wei Zhu; Zhang Chen; Chao Pei (2023) \href{https://ieeexplore.ieee.org/document/10364112/authors\#authors}{\textcolor{red}{Financial Time Series Prediction via Neural Ordinary Differential Equations Approach}} International Annual Conference on Complex Systems and Intelligent Science
        \item Hongyuan Mei, Jason Eisner (2017) \href{https://arxiv.org/pdf/1612.09328}{\textcolor{red}{The Neural Hawkes Process: A Neurally Self-Modulating Multivariate Point Process}}, NIPS
        \item Junteng Jia, Austin R. Benson (2019) \href{https://arxiv.org/pdf/1905.10403}{\textcolor{red}{Neural Jump Stochastic Differential Equations}}, NeurIPS
        \item Emmanuel Bacry, Iacopo Mastromatteo, Jean-Francois Muzy (2015) \href{https://arxiv.org/pdf/1502.04592}{\textcolor{red}{Hawkes processes in finance}}, Market Microstructure and Liquidity, 1(1).
        \item Justin Sirignano, Apaar Sadhwani, Kay Giesecke (2021) \href{https://arxiv.org/pdf/1607.02470}{\textcolor{red}{Deep Learning for Mortgage Risk}}, Journal of Financial Econometrics, 19(2), 313--368.
      \end{itemize}
  \end{itemize}

%% ---------------------------------------------------------------- Lecture 7
\item \textbf{Lecture 7: Generative diffusion Modeling for Financial Synthetic Data}
  \begin{itemize}
    \item \textbf{Scope:} This lecture introduces the generative diffusion models,
      provides and overview of its primary components and typical architecture, and
      derive the underlying equations for training. In addition, it illustrates its
      application to synthetic financial data generation.
    \item \textbf{Readings:}
      \begin{itemize}
        \item Chapter 20, Diffusion Models, from the primary textbook.
        \item Stanley Chan , (2025) \href{https://arxiv.org/pdf/2403.18103}{\textcolor{red}{Tutorial on Diffusion Models for Imaging and Vision}}
        \item Andrew Lesniewski and Giulio Trigila (2025) \href{https://arxiv.org/pdf/2412.00036}{\textcolor{red}{Beyond Monte Carlo: Harnessing Diffusion Models to Simulate Financial Market Dynamics}}
        \item Oualid Missaoui (2025) \href{https://papers.ssrn.com/sol3/papers.cfm?abstract_id=5517384}{\textcolor{red}{Deep Generative Modeling for Covariance Denoising in Finance}}
        \item Zhuohan Wang, Carmine Ventre (2024) \href{https://kclpure.kcl.ac.uk/ws/portalfiles/portal/313573729/A_Financial_Time_Series_Denoiser_Based_on_Diffusion_Model_ICAIF2024.pdf}{\textcolor{red}{A Financial Time Series Denoiser Based on Diffusion Models}}, ICAIF
      \end{itemize}
  \end{itemize}

%% ================================================================
%% Lecture 8: Graph Neural Networks for Financial Data -- DROPPED.
%% Cut for Fall 2026: the course is capped at seven sessions.
%% Kept commented out below so it can be restored if a slot opens.
%% ================================================================
% %% ---------------------------------------------------------------- Lecture 8
% \item \textbf{Lecture 8: Graph Neural Networks for Financial Data}
%   \begin{itemize}
%     \item \textbf{Scope:} This lecture will introduce Graph Neural Networks (GNNs) and
%       explain GNN architecture as well as graph embeddings. It will also illustrate
%       how can GNN be potentially leveraged to model interconnected assets or networks
%       of trading relationships.
%     \item \textbf{Readings:}
%       \begin{itemize}
%         \item Chapter 13, Graph Neural Networks, from the primary textbook.
%         \item Franco Scarselli; Marco Gori; Ah Chung Tsoi; Markus Hagenbuchner; Gabriele Monfardini (2009), \href{https://ieeexplore.ieee.org/document/4700287}{\textcolor{red}{The Graph Neural Network Model}} IEEE Transactions on Neural Networks
%         \item Gregoire Pacreau, Edmond Lezmi, Jiali Xu, (2022) \href{https://papers.ssrn.com/sol3/papers.cfm?abstract_id=3976168}{\textcolor{red}{Graph Neural Networks for Asset Management}}
%       \end{itemize}
%   \end{itemize}
\end{itemize}

\end{document}
